% !Mode:: "TeX:UTF-8"

\hitsetup{
  %******************************
  % 注意：
  %   1. 配置里面不要出现空行
  %   2. 不需要的配置信息可以删除
  %******************************
  %
  %=====
  % 秘级
  %=====
  statesecrets={公开},
  natclassifiedindex={TM301.2},
  intclassifiedindex={62-5},
  %
  %=========
  % 中文信息
  %=========
  ctitleone={局部多孔质气体静压},%本科生封面使用
  ctitletwo={轴承关键技术的研究},%本科生封面使用
  ctitlecover={局部多孔质气体静压轴承关键技术的研究},%放在封面中使用，自由断行
  ctitle={局部多孔质气体静压轴承关键技术的研究},%放在原创性声明中使用
  % csubtitle={一条副标题}, %一般情况没有，可以注释掉
  cxueke={工学},
  csubject={机械制造及其自动化},
  caffil={机电工程学院},
  cauthor={某某某},
  csupervisor={某某某教授},
  % cassosupervisor={某某某教授}, % 副指导老师
  % ccosupervisor={某某某教授}, % 联合指导老师
  % 如果是第一封面的日期要手动设置，需要取消注释下一行，并将内容改为“规范”中要求的封面第一页最下方的日期
  firstpagecdate={2022年6月},
  % 日期自动使用当前时间，若需指定按如下方式修改：
  cdate={2022年6月},
  cstudentid={2022110000},
  % cstudenttype={学术学位论文}, %非全日制教育申请学位者
  cnumber={no9527}, %编号
  cpositionname={哈铁西站}, %博士后站名称
  cfinishdate={20XX年X月---20XX年X月}, %到站日期
  csubmitdate={20XX年X月}, %出站日期
  cstartdate={3050年9月10日}, %到站日期
  cenddate={3090年10月10日}, %出站日期
  %（同等学力人员）、（工程硕士）、（工商管理硕士）、
  %（高级管理人员工商管理硕士）、（公共管理硕士）、（中职教师）、（高校教师）等
  %
  %
  %=========
  % 英文信息
  %=========
  etitle={Research on key technologies of partial porous externally pressurized gas bearing},
  % esubtitle={This is the sub title},
  exueke={Engineering},
  esubject={Computer Science and Technology},
  eaffil={\emultiline[t]{School of Mechatronics Engineering \\ Mechatronics Engineering}},
  eauthor={Yu Dongmei},
  esupervisor={Professor XXX},
  eassosupervisor={XXX},
  % 日期自动生成，若需指定按如下方式修改：
  edate={June, 2022},
  % estudenttype={Master of Art},
  %
  % 关键词用“英文逗号”分割
  ckeywords={多孔质石墨, 气体静压轴承, 稳定性},
  ekeywords={porous graphite, externally pressurized gas bearing, Stability},
}

\begin{cabstract}
气体静压轴承具有运动精度高、摩擦损耗小、发热变形小、寿命长和无污染等特点，在航空航天工业、半导体工业、纺织工业和测量仪器中得到广泛应用。本文在分析国内外气体静压轴承研究现状的基础上，以改善气体静压轴承的静态特性和稳定性为目的，通过理论分析、仿真计算和实验研究，对局部多孔质气体静压止推轴承进行了系统研究。

本文建立了基于分形几何理论的多孔质石墨渗透率与分形维数之间关系的数学模型。该模型可预测多孔质石墨的渗透率，并可直观描述孔隙大小对渗透率的影响。

在理论分析的基础上，本文建立了局部多孔质气体静压止推轴承静态特性的数学模型，并通过工程方法和有限元方法进行求解。研究结果可为局部多孔质气体静压轴承的设计和工程应用提供参考。
\end{cabstract}

\begin{eabstract}
Externally pressurized gas bearings have been widely used in aviation, semiconductor manufacturing, textile engineering and measurement instruments because of their high motion accuracy, low friction loss, low thermal deformation, long service life and clean operation. Based on a review of domestic and international studies, this thesis investigates partial porous externally pressurized gas thrust bearings through theoretical analysis, numerical simulation and experimental research.

Based on fractal geometry, a mathematical model is established to describe the relationship between the permeability of porous graphite and the fractal dimension. The model can predict the permeability of porous graphite and describe the influence of pore size on permeability.

On this basis, a mathematical model for the static characteristics of partial porous externally pressurized gas thrust bearings is developed and solved using engineering methods and the finite element method. The results provide a reference for the design and engineering application of partial porous externally pressurized gas bearings.
\end{eabstract}
